% GENERATED FILE. Edit canonical lessons, then run npm run generate:books.
% canonical-source-hash: fnv1a64:49a4f134abab5a1c
% canonical-lessons: JA-C423-gochisou, JA-C423-kanjou, JA-C423-warikan, JA-C423-ogoru, JA-C423-manpuku

\chapter{A Treat and the Bill}
\label{ch:ja-a-treat-and-the-bill}
\hlchaptermodality{423}

\begin{chapteropening}
\textbf{By the end of this chapter:} \emph{I can treat a friend and split the bill.}

Complete the last lesson of chapter 423: I can treat a friend and split the bill.
\end{chapteropening}

\section[\texorpdfstring{\ja{ごちそう} (gochisō)}{ごちそう (gochisō)}]{\ja{ごちそう} (gochisō) — a feast; a treat}

\label{lesson:JA-C423-gochisou}



\begin{quote}
Before the new one: say the Japanese for sliced raw fish, then the Japanese for grilled chicken on skewers.
\end{quote}

\subsection*{You'll want to know: \ja{ごちそう}}
\textbf{\ja{ごちそう}} — \emph{gochisō} — \textquotedblleft{}a feast; a treat\textquotedblright{}.

\begin{quote}
\textbf{\ja{そふの} \ja{うちで} \ja{ごちそうを} \ja{たべました。}} — \emph{sofu no uchi de gochisō o tabemashita} — \textquotedblleft{}We had a feast at my grandfather's house.\textquotedblright{}
\end{quote}

\begin{cousinweb}
\textbf{\ja{ごちそう}} (\emph{gochisō}) is Sino-Japanese (kango): the polite prefix \emph{go-} (Middle Chinese \emph{ngjoH}) and two Chinese characters read with their on'yomi, readings borrowed from Chinese: \emph{chi}, \textquotedblleft{}gallop\textquotedblright{} (\emph{drje}), and \emph{sō}, \textquotedblleft{}run\textquotedblright{} (\emph{tsuwX}); JMdict gives the sense \textquotedblleft{}a treat\textquotedblright{} without saying how the running came to mean it (KANJIDIC2; Middle Chinese in Baxter's transcription).
\end{cousinweb}

\begin{grammarlens}[title={What you've built}]
One more word for getting around.
\end{grammarlens}

\subsection*{Your turn}
\begin{itemize}
\raggedright
  \item \emph{Say it:} \emph{gochisō}
  \item \emph{Say it:} \emph{gochisō}, once more
  \item \emph{Say it:} \emph{sofu no uchi de gochisō o tabemashita} — We had a feast at my grandfather's house.
  \item \emph{Recall:} say the Japanese for sliced raw fish, then the Japanese for grilled chicken on skewers, then say \emph{gochisō} again
\end{itemize}

\begin{tcolorbox}[breakable,colback=teal!4,colframe=teal!35!black,title={Before you move on}]
What does \ja{ごちそう} mean? (\textquotedblleft{}A feast; a treat\textquotedblright{}.) Say it once more.
\end{tcolorbox}
\section[\texorpdfstring{\ja{かんじょう} (kanjō)}{かんじょう (kanjō)}]{\ja{かんじょう} (kanjō) — paying the bill; a calculation}

\label{lesson:JA-C423-kanjou}



\begin{quote}
Before the new one: say the Japanese for grilled chicken on skewers, then the Japanese for a feast.
\end{quote}

\subsection*{You'll want to know: \ja{かんじょう}}
\textbf{\ja{かんじょう}} — \emph{kanjō} — \textquotedblleft{}paying the bill; a calculation\textquotedblright{}.

\begin{quote}
\textbf{\ja{かんじょうを} \ja{おねがいします。}} — \emph{kanjō o onegai shimasu} — \textquotedblleft{}The bill, please.\textquotedblright{}
\end{quote}

\begin{cousinweb}
\textbf{\ja{かんじょう}} (\emph{kanjō}) is Sino-Japanese (kango): two Chinese characters read with their on'yomi, readings borrowed from Chinese. They are \emph{kan}, \textquotedblleft{}check\textquotedblright{} (Middle Chinese \emph{khomH}) and \emph{jō}, \textquotedblleft{}determine\textquotedblright{} (\emph{dengH}) (KANJIDIC2; Middle Chinese in Baxter's transcription).
\end{cousinweb}

\begin{grammarlens}[title={What you've built}]
One more word for getting around.
\end{grammarlens}

\subsection*{Your turn}
\begin{itemize}
\raggedright
  \item \emph{Say it:} \emph{kanjō}
  \item \emph{Say it:} \emph{kanjō}, once more
  \item \emph{Say it:} \emph{kanjō o onegai shimasu} — The bill, please.
  \item \emph{Recall:} say the Japanese for grilled chicken on skewers, then the Japanese for a feast, then say \emph{kanjō} again
\end{itemize}

\begin{tcolorbox}[breakable,colback=teal!4,colframe=teal!35!black,title={Before you move on}]
What does \ja{かんじょう} mean? (\textquotedblleft{}Paying the bill; a calculation\textquotedblright{}.) Say it once more.
\end{tcolorbox}
\section[\texorpdfstring{\ja{わりかん} (warikan)}{わりかん (warikan)}]{\ja{わりかん} (warikan) — splitting the bill}

\label{lesson:JA-C423-warikan}



\begin{quote}
Before the new one: say the Japanese for a feast, then the Japanese for paying the bill.
\end{quote}

\subsection*{You'll want to know: \ja{わりかん}}
\textbf{\ja{わりかん}} — \emph{warikan} — \textquotedblleft{}splitting the bill\textquotedblright{}.

\begin{quote}
\textbf{\ja{ともだちと} \ja{わりかんで} \ja{はらいました。}} — \emph{tomodachi to warikan de haraimashita} — \textquotedblleft{}My friend and I split the bill.\textquotedblright{}
\end{quote}

\begin{cousinweb}
\textbf{\ja{わりかん}} (\emph{warikan}) is built from both layers of the vocabulary: the native Japanese (wago) \emph{wari}, \textquotedblleft{}splitting\textquotedblright{} (from \emph{waru}), inherited rather than borrowed, and the Sino-Japanese \emph{kan}, \textquotedblleft{}check\textquotedblright{} (Middle Chinese \emph{khomH}), an on'yomi borrowed from Chinese (JMdict; KANJIDIC2; Middle Chinese in Baxter's transcription).
\end{cousinweb}

\begin{grammarlens}[title={What you've built}]
One more word for getting around.
\end{grammarlens}

\subsection*{Your turn}
\begin{itemize}
\raggedright
  \item \emph{Say it:} \emph{warikan}
  \item \emph{Say it:} \emph{warikan}, once more
  \item \emph{Say it:} \emph{tomodachi to warikan de haraimashita} — My friend and I split the bill.
  \item \emph{Recall:} say the Japanese for a feast, then the Japanese for paying the bill, then say \emph{warikan} again
\end{itemize}

\begin{tcolorbox}[breakable,colback=teal!4,colframe=teal!35!black,title={Before you move on}]
What does \ja{わりかん} mean? (\textquotedblleft{}Splitting the bill\textquotedblright{}.) Say it once more.
\end{tcolorbox}
\section[\texorpdfstring{\ja{おごる} (ogoru)}{おごる (ogoru)}]{\ja{おごる} (ogoru) — to treat someone; to buy someone a meal}

\label{lesson:JA-C423-ogoru}



\begin{quote}
Before the new one: say the Japanese for paying the bill, then the Japanese for splitting the bill.
\end{quote}

\subsection*{You'll want to know: \ja{おごる}}
\textbf{\ja{おごる}} — \emph{ogoru} — \textquotedblleft{}to treat someone; to buy someone a meal\textquotedblright{}.

\begin{quote}
\textbf{\ja{ともだちに} \ja{ひるごはんを} \ja{おごりました。}} — \emph{tomodachi ni hirugohan o ogorimashita} — \textquotedblleft{}I treated my friend to lunch.\textquotedblright{}
\end{quote}

Polite \textbf{\ja{おごります}}, \emph{ogorimasu}; a godan verb. Past (polite): \textbf{\ja{あにが} \ja{ばんごはんを} \ja{おごりました。}}, \emph{ani ga bangohan o ogorimashita}.

\begin{cousinweb}
\textbf{\ja{おごる}} (\emph{ogoru}) is a native Japanese word (wago), inherited rather than borrowed from Chinese; its kanji was attached later as a kun'yomi (JMdict).
\end{cousinweb}

\begin{grammarlens}[title={What you've built}]
One more word for getting around.
\end{grammarlens}

\subsection*{Your turn}
\begin{itemize}
\raggedright
  \item \emph{Say it:} \emph{ogoru}
  \item \emph{Say it:} \emph{ogoru}, once more
  \item \emph{Say it:} \emph{tomodachi ni hirugohan o ogorimashita} — I treated my friend to lunch.
  \item \emph{Recall:} say the Japanese for paying the bill, then the Japanese for splitting the bill, then say \emph{ogoru} again
\end{itemize}

\begin{tcolorbox}[breakable,colback=teal!4,colframe=teal!35!black,title={Before you move on}]
What does \ja{おごる} mean? (\textquotedblleft{}To treat someone; to buy someone a meal\textquotedblright{}.) Say it once more.
\end{tcolorbox}
\section[\texorpdfstring{\ja{まんぷく} (manpuku)}{まんぷく (manpuku)}]{\ja{まんぷく} (manpuku) — a full stomach}

\label{lesson:JA-C423-manpuku}



\begin{quote}
Before the new one: say the Japanese for splitting the bill, then the Japanese for to treat someone.
\end{quote}

\subsection*{You'll want to know: \ja{まんぷく}}
\textbf{\ja{まんぷく}} — \emph{manpuku} — \textquotedblleft{}a full stomach\textquotedblright{}.

\begin{quote}
\textbf{\ja{ごちそうを} \ja{たべて、まんぷくに} \ja{なりました。}} — \emph{gochisō o tabete, manpuku ni narimashita} — \textquotedblleft{}I ate the feast and was full.\textquotedblright{}
\end{quote}

\begin{cousinweb}
\textbf{\ja{まんぷく}} (\emph{manpuku}) is Sino-Japanese (kango): two Chinese characters read with their on'yomi, readings borrowed from Chinese. They are \emph{man}, \textquotedblleft{}full\textquotedblright{} (Middle Chinese \emph{manX}) and \emph{fuku}, \textquotedblleft{}belly\textquotedblright{} (\emph{pjuwk}), joined as \emph{manpuku} (KANJIDIC2; Middle Chinese in Baxter's transcription).
\end{cousinweb}

\begin{grammarlens}[title={What you've built}]
One more word for getting around.
\end{grammarlens}

\subsection*{Your turn}
\begin{itemize}
\raggedright
  \item \emph{Say it:} \emph{manpuku}
  \item \emph{Say it:} \emph{manpuku}, once more
  \item \emph{Say it:} \emph{gochisō o tabete, manpuku ni narimashita} — I ate the feast and was full.
  \item \emph{Recall:} say the Japanese for splitting the bill, then the Japanese for to treat someone, then say \emph{manpuku} again
\end{itemize}

\begin{tcolorbox}[breakable,colback=teal!4,colframe=teal!35!black,title={Before you move on}]
What does \ja{まんぷく} mean? (\textquotedblleft{}A full stomach\textquotedblright{}.) Say it once more.
\end{tcolorbox}
